\chapter{Assets and Followers}
łabel{ch:assets−followers}
∈dex{assets}∈dex{followers}

Power in Fate's Edge is not measured only by what your character can do with their own hands.  
As characters grow, their influence extends outward−−−into places, people, and institutions that act in their name.

\textbf{Assets} and \textbf{Followers} represent this outward reach.  
They are how a character's past choices continue to matter when the dice stop rolling.

\emph{Assets reshape the world around you.}  
\emph{Followers stand beside you when the world pushes back.}

\begin{tcolorbox}[colback=green!5!white,colframe=green!75!black,title=A Word from the Raven Road]
\emph{``A safehouse is not a home; it is a promise you may not keep. A follower is not a servant; they are a weight you choose to carry. Invest wisely, or the road will collect.''}  
\vspace{0.2cm}ℎfill −−− Dusana of the Raven Road, \textit{The Hearth Ledger}
\end{tcolorbox}

\section{Understanding Assets and Followers}
łabel{sec:assets−followers−overview}
∈dex{assets!overview}∈dex{followers!overview}

Assets and Followers serve different roles in play, but both are expressions of legacy, leverage, and risk.

\subsection*{Key Differences}
\begin{itemize}
ℹtem \textbf{Assets} are off‑screen resources: places, networks, titles, or enterprises.
ℹtem \textbf{Followers} are on‑screen allies who act during scenes.
ℹtem \textbf{Assets} change the situation \emph{before or around} a scene.
ℹtem \textbf{Followers} participate \emph{directly} in the moment.
\end{itemize}

Both are powerful−−−but neither is free.

\subsection*{Management and Risk}
Assets and Followers require care and invite consequences:
\begin{itemize}
ℹtem They demand upkeep through XP, downtime, or narrative obligation.
ℹtem They can be targeted, compromised, or destroyed.
ℹtem Overuse draws attention from enemies, factions, and rivals.
\end{itemize}

\emph{Power creates gravity. What you rely on can be taken from you.}

\section{Assets System}
łabel{sec:assets−system}
∈dex{assets!system}

Assets are persistent resources under your control: holdings, institutions, reputations, or infrastructures.

\textit{\textbf{Assets do not replace rolls. They change what rolls are possible.}}

Assets allow characters to:
\begin{itemize}
ℹtem Arrive prepared instead of improvising.
ℹtem Reframe danger into opportunity.
ℹtem Shift Position without negating risk.
\end{itemize}

\subsection*{Asset Types and Costs}
łabel{sec:asset−types−costs}
\begin{center}
\begin{tabular}{lll}
⊤rule
\textbf{Type} & \textbf{XP Cost} & \textbf{Establishment Time} \\
∣rule
Minor    & 4 XP  & 1 day \\
Standard & 8 XP  & 1 week \\
Major    & 12 XP & 1 month \\
⊥tomrule
\end{tabular}
\end{center}

\subsection*{Asset Examples}
łabel{sec:asset−examples}
\begin{description}
ℹtem[Minor Assets] Safehouse, small shop, minor title, basic workshop.
ℹtem[Standard Assets] Guild membership, noble office, trading post, spy network.
ℹtem[Major Assets] Fortress, chartered city license, major enterprise, regional authority.
\end{description}

\subsection*{Using Assets}
łabel{sec:using−assets}

Assets provide leverage in two primary ways: quiet preparation and dramatic activation.

\paragraph{Free Off‑Screen Use}
łabel{sec:asset−free−use}
Each Asset has a defined off‑screen effect usable \textbf{once per session}, typically during downtime or scene framing:
\begin{itemize}
ℹtem \textbf{Safehouse}: Secure lodging or concealment.
ℹtem \textbf{Spy Network}: Baseline intelligence about a person or place.
ℹtem \textbf{Workshop}: Repair gear or prepare simple solutions.
ℹtem \textbf{Trading Post}: Acquire common goods or establish contacts.
\end{itemize}
These uses prepare the ground−−−they do not resolve active conflict.

\paragraph{Boon Activation}
łabel{sec:asset−boon−activation}
Spend \textbf{1 Boon} (see \ref{ch:core−mechanics}) to activate an Asset dramatically during a scene.

This does not guarantee success. Instead, the Asset \textbf{reshapes the fiction} so a different outcome becomes possible.

An activated Asset may:
\begin{itemize}
  ℹtem Improve \textbf{Position} by one step.
  ℹtem Reduce the effective \textbf{DV} of a follow‑up action.
  ℹtem Advance, suppress, or redirect a relevant Timer .
  ℹtem Grant temporary access, protection, or narrative permission.
\end{itemize}

\textbf{Examples:}
\begin{itemize}
  ℹtem \textbf{Safehouse}: Reveal a concealed exit, allowing escape at \emph{Controlled} Position.
  ℹtem \textbf{Spy Network}: Surface critical intelligence, advancing an investigation Timer .
  ℹtem \textbf{Workshop}: Jury‑rig a solution that enables an otherwise impossible attempt.
  ℹtem \textbf{Trading Post}: Call in a favor to gain leverage or access in a negotiation.
\end{itemize}

The GM determines the precise effect based on the Asset's scale, condition, and the current stakes.

\paragraph{High‑Tier Use}
łabel{sec:asset−high−tier}
At higher tiers (see \ref{sec:tiers}), Assets carry greater narrative weight−−−and greater risk.

Repeated or aggressive use may:
\begin{itemize}
  ℹtem Create new complications.
  ℹtem Advance hostile faction Timers.
  ℹtem Place the Asset itself in danger.
\end{itemize}
\emph{Assets are leverage, not immunity.}

\paragraph{XP Activation}
łabel{sec:asset−xp−activation}
Spend \textbf{2 XP} (see \ref{ch:advancement}) to use an Asset's off‑screen effect beyond its normal limit:
\begin{itemize}
ℹtem Emergency use after the free activation is spent.
ℹtem Additional downtime applications.
ℹtem Exceptional circumstances requiring extended support.
\end{itemize}
XP use represents strain, overreach, or long‑term cost.

\section{Asset Status}
łabel{sec:asset−status}
∈dex{assets!status}

Assets have a single unified \textbf{Status} that describes their overall health and usability.

\begin{center}
\small
\begin{tabular}{lll}
⊤rule
\textbf{Status} & \textbf{Meaning} & \textbf{Effect} \\
∣rule
Flourishing & Well−maintained, respected & +1 die on rolls involving the asset, +1 Yield \\
Stable & Normal operation & No modifier \\
Strained & Problems accumulating & –1 die on asset−related rolls, no Yield \\
Collapsed & Non−functional & Requires quest or major downtime to restore \\
⊥tomrule
\end{tabular}
\end{center}

The Status worsens when the asset is neglected, attacked, or suffers a public failure. It improves through dedicated downtime actions or successful missions that benefit the asset.

\subsection*{Maintenance Requirements}
łabel{sec:asset−maintenance}
\begin{itemize}
ℹtem \textbf{Regular Upkeep}: Two options (see \ref{subsec:upkeep−options}).
ℹtem \textbf{Neglect}: Assets deteriorate if not maintained (Status worsens one step).
ℹtem \textbf{Recovery}: Collapsed assets require significant effort to restore, typically a quest or a full downtime arc.
\end{itemize}

\section{Followers System}
łabel{sec:followers−system}
∈dex{followers!system}

Followers are characters who assist you directly.

\subsection*{Follower Capability Ratings}
łabel{sec:follower−cap}
Followers are rated by Capability (\textbf{Cap}) from 1 to 5:

\begin{center}
\begin{tabular}{ll}
⊤rule
\textbf{Cap} & \textbf{Description} \\
∣rule
1 & Novice helper, basic assistance \\
2 & Competent assistant, reliable support \\
3 & Skilled specialist, valuable aid \\
4 & Expert ally, significant capability \\
5 & Master companion, exceptional ability \\
⊥tomrule
\end{tabular}
\end{center}

\subsection*{Follower Costs}
łabel{sec:follower−costs}
\begin{itemize}
ℹtem \textbf{XP Cost}: Capability squared ($\text{Cap}²$).
ℹtem \textbf{Example}: Cap 3 follower costs $3² = 9$ XP.
ℹtem \textbf{Recruitment}: 1−−3 days downtime to find and brief.
ℹtem \textbf{Limits}: The GM may set maximum followers based on story.
\end{itemize}

\subsection*{Follower Types}
łabel{sec:follower−types}
\begin{description}
ℹtem[Combat Allies] Warriors, guards, mercenaries.
ℹtem[Technical Experts] Craftspeople, engineers, specialists.
ℹtem[Social Contacts] Informants, diplomats, agents.
ℹtem[Specialists] Unique capabilities like magic or stealth.
\end{description}

\subsection*{Follower States}
łabel{sec:follower−states}
Instead of numeric timers, track a follower's \textbf{Loyalty} and \textbf{Fitness} as one of three states each.

\begin{center}
\small
\begin{tabular}{lll}
⊤rule
\textbf{State} & \textbf{Meaning} & \textbf{Mechanical Effect} \\
∣rule
μlticolumn{3}{c}{\textit{Loyalty}} \\
Faithful & Committed, trustworthy & Normal assist, independent actions available \\
Strained & Doubtful, pressured & Assist at –1 die; may refuse risky orders \\
Broken & Betrayed or deserted & Follower leaves or becomes hostile \\
∣rule
μlticolumn{3}{c}{\textit{Fitness}} \\
Ready & Uninjured, rested & Full function \\
Hurt & Wounded, exhausted & Assist at –1 die; no independent actions \\
Down & Incapacitated & Requires rescue or medical care \\
⊥tomrule
\end{tabular}
\end{center}

Change states when the fiction demands it (e.g., after a failed mission, a moral compromise, or taking damage). No rolls required.

\subsection*{High‑Capability Followers (Cap 4–5)}
łabel{sec:high−cap−followers}
Followers with Capability 4 or 5 are trusted lieutenants, stewards, or commanders. They can be delegated the management of your Assets, freeing you from direct oversight.

\begin{itemize}
  ℹtem A \textbf{Cap 4} follower can be delegated \textbf{up to 2 Assets}.
  ℹtem A \textbf{Cap 5} follower can be delegated \textbf{up to 4 Assets}.
\end{itemize}

When an Asset is delegated to a follower, the follower's own upkeep (paid in XP or downtime action) \textbf{automatically covers the upkeep of those Assets} — you do not pay separate Asset upkeep. The follower handles the logistics, staff, or political maintenance as part of their role.

If the follower's Loyalty becomes \textbf{Strained} or \textbf{Broken}, all delegated Assets immediately worsen one step (e.g., Flourishing $→$ Stable, Stable $→$ Strained) until the follower is restored. If a follower becomes \textbf{Broken} (leaves or betrays), those Assets are lost.

\section{Using Followers}
łabel{sec:using−followers}
∈dex{followers!usage}

\subsection*{Assistance in Scenes}
łabel{sec:follower−assistance}
Followers can help with your actions:
\begin{itemize}
ℹtem \textbf{Assist Dice}: Add dice equal to $\min(\text{Cap}, \text{relevant skill})$.
ℹtem \textbf{Maximum Bonus}: +3 dice total from all sources (see \ref{ch:core−mechanics}).
ℹtem \textbf{Cost}: Spend 1 Boon or 1 Fatigue to add +1 die (max +3 from assists).
ℹtem \textbf{One Helper}: Only one follower can assist per action.
\end{itemize}

\subsection*{Independent Actions}
łabel{sec:follower−independent}
Once per scene (party‑wide), a follower can take a small action:
\begin{description}
ℹtem[\emph{Scout & Signal}] Change an ally's next action to \textbf{Dominant} position.
ℹtem[\emph{Distract & Draw}] Reduce a threat timer by 1 segment.
ℹtem[\emph{Fetch & Carry}] Move an object through danger safely.
\end{description}

\subsection*{Cost of Independent Actions}
łabel{sec:follower−independent−cost}
\begin{itemize}
ℹtem Worsen \textbf{Fitness} one step (Ready $→$ Hurt $→$ Down), \emph{or}
ℹtem Worsen \textbf{Loyalty} one step (Faithful $→$ Strained $→$ Broken).
ℹtem Cannot be used if the follower is already \textbf{Broken} or \textbf{Down}.
\end{itemize}

\subsection*{Off‑Screen Capability}
łabel{sec:follower−offscreen}
Once per downtime, a \textbf{Cap 5} follower can solve a significant problem:
\begin{itemize}
ℹtem Generates \textbf{1 Story Beat} for the GM (the GM describes a resulting complication).
ℹtem Useful for emergencies, but costly to the story.
\end{itemize}

\section{Upkeep and Maintenance}
łabel{subsec:upkeep−options}
∈dex{upkeep}∈dex{assets!upkeep}∈dex{followers!upkeep}

Both assets and followers require regular maintenance. Each downtime period, choose one of the two options.

\subsection*{Upkeep Options}
łabel{sec:upkeep−options}
\begin{itemize}
ℹtem \textbf{Efficient (Higher XP, Less Time)}:  
  Pay Upkeep XP = $\max(1, ⌈ \text{Acquisition XP} / 3 ⌉)$, minimal effort.  
  \emph{Rule of thumb: about one‑third the acquisition cost, minimum 1 XP.}
ℹtem \textbf{Intensive (Lower XP, More Time)}:  
  Pay 1 XP, use a dedicated downtime action with significant personal involvement.
ℹtem \textbf{Failure to Pay}:  
  Asset Status worsens one step (Flourishing $→$ Stable $→$ Strained $→$ Collapsed).  
  Follower Loyalty worsens one step (Faithful $→$ Strained $→$ Broken).
\end{itemize}

\subsection*{Upkeep Examples}
łabel{sec:upkeep−examples}
\begin{itemize}
  ℹtem \textbf{Minor Asset (4 XP)}: Efficient = 2 XP; Intensive = 1 XP + action.
  ℹtem \textbf{Standard Asset (8 XP)}: Efficient = 3 XP; Intensive = 1 XP + action.
  ℹtem \textbf{Cap 3 Follower (9 XP)}: Efficient = 3 XP; Intensive = 1 XP + action.
  ℹtem \textbf{Cap 2 Follower (4 XP)}: Efficient = 2 XP; Intensive = 1 XP + action.
\end{itemize}

\section{Strategic Considerations}
łabel{sec:asset−follower−strategy}
∈dex{strategic considerations}

\subsection*{When to Invest in Assets}
łabel{sec:invest−assets}
\begin{itemize}
ℹtem You need reliable off‑screen capabilities.
ℹtem Your character concept involves wealth or influence.
ℹtem The party lacks certain logistical support.
ℹtem You want to build long‑term influence.
\end{itemize}

\subsection*{When to Invest in Followers}
łabel{sec:invest−followers}
\begin{itemize}
ℹtem You need on‑screen assistance.
ℹtem Your character works better with support.
ℹtem The party needs specific capabilities you lack.
ℹtem You want character‑driven story opportunities.
\end{itemize}

\subsection*{Investment by Player Archetype}
łabel{sec:invest−by−archetype}
\begin{center}
\begin{tabular}{lll}
⊤rule
\textbf{Archetype} & \textbf{Typical XP Investment} & \textbf{Best For} \\
∣rule
Solo & 0−−10% & Reliable personal power, minimal upkeep. \\
Mixed (Balanced) & 15−−░ & One solid Follower or Asset (scout, safehouse). \\
Mastermind & 35−−55% & Networks, informants, multiple Assets; build a faction. \\
⊥tomrule
\end{tabular}
\end{center}

\subsection*{Balance Recommendations}
łabel{sec:asset−follower−balance}
\begin{itemize}
ℹtem \textbf{Personal Path} (see \ref{sec:xp−personal−path}): 0−−10% of XP in assets/followers.
ℹtem \textbf{Balanced Path} (see \ref{sec:xp−balanced−path}): 15−−░.
ℹtem \textbf{Influencer Path} (see \ref{sec:xp−influencer−path}): 35−−55%.
\end{itemize}

\begin{tcolorbox}[colback=red!5!white,colframe=red!75!black,title=Planning for Succession (The Legacy Engine)]
∈dex{succession}∈dex{legacy engine}
A loyal Follower (Cap 3+) can become your \textbf{successor} when your character retires or falls −− inheriting your faction standing, one unresolved burden, and a Legacy Bond. See \textbf{Chapter \ref{ch:legacy−engine}: The Legacy Engine} for full rules. If you plan for your character's story to continue beyond them, invest in a follower you would trust to carry your banner.
\end{tcolorbox}

\section{Terrestrial Patrons}
łabel{sec:terrestrial−patrons}
∈dex{Patrons!terrestrial}

Not all power comes from cosmic entities or divine forces. Sometimes your greatest allies−−−and most demanding creditors−−−are mortal factions who can offer protection, resources, and influence in the world of men.

\subsection{What Are Terrestrial Patrons?}
łabel{sec:terrestrial−definition}

Terrestrial Patrons represent your ongoing relationships with powerful mortal organizations:
\begin{itemize}
    ℹtem Noble houses and courtly factions.
    ℹtem Merchant guilds and trading consortiums.
    ℹtem Military units and mercenary companies.
    ℹtem Criminal organizations and smuggling rings.
    ℹtem Religious orders and temple hierarchies.
    ℹtem Scholarly colleges and arcane societies.
\end{itemize}

Unlike supernatural Patrons (see \ref{ch:magic}), Terrestrial Patrons do not grant magic, but they offer \textbf{leverage}: protection, resources, sanctuary, information, and political influence.

\subsection{Gaining a Terrestrial Patron}
łabel{sec:terrestrial−gaining}

To gain a Terrestrial Patron, complete one of these significant actions:
\begin{itemize}
    ℹtem Complete a major job or service for them.
    ℹtem Swear a formal Oath of service.
    ℹtem Enter into legal or financial binding.
    ℹtem Share criminal secrets or be compromised by them.
    ℹtem Perform a notable service that earns their favor.
\end{itemize}
Mark the Patron on your character sheet and write one sentence: \emph{``They want me because _______''}.

\subsection{Obligation System}
łabel{sec:terrestrial−obligation}

Terrestrial Patrons use the same Obligation track as supernatural Patrons (see \ref{sec:obligation−capacity}), but consequences are social, legal, or economic rather than mystical.

\begin{itemize}
    ℹtem \textbf{Calling in Favors}: When you request Patron influence, add +1 Obligation.
    ℹtem \textbf{Capacity}: Spirit + Presence (same as supernatural Patrons).
    ℹtem \textbf{Overflow}: Each segment above capacity inflicts 1 Fatigue.
    ℹtem \textbf{Resolution}: Reduce through service, downtime actions, or fulfilling demands.
\end{itemize}

\subsection{Patron Perks}
łabel{sec:terrestrial−perks}

Each Terrestrial Patron offers 2−−3 repeatable benefits that require no rolls to access.

\textbf{Common Perks Include:}
\begin{itemize}
    ℹtem Sanctuary in their territory or establishments.
    ℹtem Legal relief or protection from certain authorities.
    ℹtem Access to black market goods or specialized services.
    ℹtem Elite followers or hirelings at reduced cost.
    ℹtem Forged documents or identity protection.
    ℹtem Information networks and rumor access.
    ℹtem Trade advantages or exclusive contracts.
\end{itemize}
Using a perk never requires a roll−−−Fate has already been paid. Each use comes with the ongoing cost of Obligation.

\subsection{Patron Demands}
łabel{sec:terrestrial−demands}

Terrestrial Patrons always want something in return:
\begin{itemize}
    ℹtem Silence about sensitive information.
    ℹtem Loyalty over other conflicting interests.
    ℹtem Completion of specific jobs or tasks.
    ℹtem Delivery of valuable names or secrets.
    ℹtem Political or social support for their causes.
    ℹtem Elimination of rivals or threats to their interests.
\end{itemize}
\textbf{Refusing a Demand} raises Obligation by 1 segment and may trigger immediate complications.

\subsection{When Obligation Fills}
łabel{sec:terrestrial−obligation−fill}

At \textbf{6 Obligation segments}, the Patron acts−−−this is not optional. Choose one:
\begin{enumerate}
    ℹtem \textbf{Unavoidable Job}: You must complete a task you cannot refuse.
    ℹtem \textbf{Severe Price}: Pay significant cost (legal, social, or material).
    ℹtem \textbf{First Strike}: They act against you−−−reputation damage, warrants, bounty, blackmail.
\end{enumerate}
After the consequence lands, reduce Obligation to 3 segments.

\subsection{Cutting Ties}
łabel{sec:terrestrial−cutting−ties}

You may sever a Terrestrial Patron relationship, but doing so has serious fallout:
\begin{itemize}
    ℹtem Lose all current perks immediately.
    ℹtem Gain a new \textbf{Rival} faction that opposes your interests.
    ℹtem Take a lasting \textbf{Complication} (Curse, Bounty, or Scandal) that follows you.
    ℹtem May trigger immediate retaliation depending on the Patron's nature.
\end{itemize}
Some Patrons never forgive betrayal; others can be bought off with significant compensation.

\subsection{Example: The Black Ledger}
łabel{sec:terrestrial−example}

The Black Ledger is a smuggling syndicate operating in the coastal city of Kahfagia.

\textbf{Gaining Their Favor:} Successfully complete three smuggling runs without detection.

\textbf{Perks:}
\begin{itemize}
    ℹtem Sanctuary in hidden safehouses throughout the city.
    ℹtem Access to rare imported goods and contraband.
    ℹtem Information about port authorities and customs schedules.
\end{itemize}

\textbf{Demands:}
\begin{itemize}
    ℹtem Maintain silence about syndicate operations.
    ℹtem Transport occasional ``special packages'' (may be illegal).
    ℹtem Eliminate rival smugglers who threaten their routes.
\end{itemize}

\textbf{Sample Scenario:} Rellan calls on the Ledger for a smuggled border crossing. The GM rules it succeeds automatically, but adds +1 Obligation. Later, the Ledger demands he silence a witness who saw too much. If he refuses, Obligation rises again. If Obligation ever reaches 6, the Ledger collects: accounts frozen, bounty posted, or a rival informant sent after him.

\section{Multi‑Character Switching (Optional)}
łabel{sec:multi−character−switching}
∈dex{Echo Character}

You may choose to temporarily step away from your current character and play another, while your first character continues to act offscreen. This is called an \textbf{Echo Character}. It allows you to explore different parts of the story without permanently retiring anyone.

\subsection{How It Works}
When you switch, you declare:
\begin{itemize}
    ℹtem Which character goes offscreen.
    ℹtem Where they are going (fictional location).
    ℹtem Why they are stepping out (story reason).
    ℹtem How long you expect to be away (1−−3 sessions).
\end{itemize}
The offscreen character remains active in the fiction. During each session they are away, you complete a simple \textbf{Downtime Log}:
\begin{itemize}
    ℹtem One \textbf{Resource Action} (how they maintain connections).
    ℹtem One \textbf{Project Action} (a goal advanced).
    ℹtem One \textbf{Event Response} (how they react to offscreen events).
\end{itemize}
The GM records the outcomes and may introduce a consequence.

\subsection{Returning to Play}
When you bring the character back:
\begin{itemize}
    ℹtem You gain \textbf{2 Boons} for narrative re‑entry.
    ℹtem Roll \textbf{Wits + Spirit} (DV 3) to reestablish presence.
    ℹtem The GM describes one consequence of your absence (e.g., an opportunity missed, a rival gained, a personal change).
\end{itemize}

For full mechanical details (Crisis Timers, Buddy System, etc.), see the GM Guide.

\section{Practical Examples}
łabel{sec:asset−follower−examples}

\subsection*{Asset Example: The Safehouse}
łabel{sec:example−safehouse}
\begin{itemize}
ℹtem \textbf{Type}: Minor Asset (4 XP)
ℹtem \textbf{Free Use}: Secure lodging, basic supplies between adventures.
ℹtem \textbf{Boon Activation}: Reveal a hidden escape route during pursuit.
ℹtem \textbf{Upkeep}: Efficient: 2 XP; Intensive: 1 XP + downtime action.
ℹtem \textbf{Risks}: Discovery by enemies, maintenance costs.
\end{itemize}

\subsection*{Follower Example: The Scout}
łabel{sec:example−scout}
\begin{itemize}
ℹtem \textbf{Capability}: 3 (9 XP)
ℹtem \textbf{Assistance}: +3 dice on tracking and survival rolls.
ℹtem \textbf{Independent Action}: Scout ahead to improve party position.
ℹtem \textbf{Upkeep}: Efficient: 3 XP; Intensive: 1 XP + action.
ℹtem \textbf{Risks}: Injury in dangerous scouting; disloyalty if mistreated.
\end{itemize}

\subsection*{Combination Example: The Merchant}
łabel{sec:example−merchant}
\begin{itemize}
ℹtem \textbf{Assets}: Trading post (8 XP), caravan (4 XP) −−− \emph{12 XP total}
ℹtem \textbf{Followers}: Cap 2 guards (4 XP each = 8 XP), Cap 3 factor (9 XP) −−− \emph{17 XP total}
ℹtem \textbf{Total Investment}: \textbf{29 XP} in assets and followers
ℹtem \textbf{Upkeep (Efficient)}: Assets 4 XP + Followers 6 XP = \textbf{10 XP} per downtime
ℹtem \textbf{Benefits}: Trade income, transport, protection, business contacts
ℹtem \textbf{Risks}: Competition, bandit attacks, employee issues, regulatory attention
\end{itemize}

\emph{Remember:} Assets and followers can greatly expand your capabilities, but they require careful management and carry significant risks. Invest wisely based on your character concept and the needs of your group.

\section{Narrative‑Heavy Options}
łabel{sec:narrative−assets}

For groups that prefer strong narrative focus, consider these optional approaches.

\paragraph{Story‑Driven Upkeep}
Instead of tracking XP costs for upkeep, the GM can introduce narrative complications that require attention. A neglected asset might attract unwanted attention, while a neglected follower might request a favor or special treatment.

\paragraph{Collaborative Management}
Describe how you maintain your assets and followers through roleplay rather than mechanical upkeep costs. A well‑described scene of tending a workshop or bonding with followers can fulfill maintenance requirements.

\paragraph{Flexible Condition Tracking}
Focus on the narrative implications of asset and follower conditions rather than strict penalties. A Strained asset might still function but with interesting complications, while a Collapsed asset might require creative solutions rather than just XP investment.

\begin{tcolorbox}[colback=blue!5!white,colframe=blue!75!black,title={Assets and Followers Quick Reference},fonttitle=\bfseries]
\textbf{Assets:}
\begin{itemize}
ℹtem Minor: 4 XP  |  Standard: 8 XP  |  Major: 12 XP
ℹtem Free off‑screen use: once per session
ℹtem Boon activation: spend 1 Boon for scene impact
ℹtem Status: Flourishing $→$ Stable $→$ Strained $→$ Collapsed
\end{itemize}

\textbf{Followers:}
\begin{itemize}
ℹtem Cost: $\text{Cap}²$ XP
ℹtem Assistance: $+\min(\text{Cap},\text{skill})$ dice (max +3 from all sources)
ℹtem Independent action: once per scene (party‑wide)
ℹtem States: Loyalty (Faithful/Strained/Broken); Fitness (Ready/Hurt/Down)
\end{itemize}

\textbf{Upkeep Options:}
\begin{itemize}
ℹtem Efficient: $\max(1,⌈\text{Cost}/3⌉)$ XP (about one‑third)
ℹtem Intensive: 1 XP + dedicated downtime action
\end{itemize}
\end{tcolorbox}